\documentclass[sigconf]{acmart}

\settopmatter{printacmref=false,printccs=false}
\renewcommand\footnotetextcopyrightpermission[1]{} % removes footnote with conference information in first column

%% CBSoft uses an ACM-like format without a copyright note.
\setcopyright{none}

%% TODO: conferir edição e data quando sair a chamada do CBSoft 2027.
\acmConference[SBES 2027]{40th Brazilian Symposium on Software Engineering}{September 27--30, 2026}{Belo Horizonte, MG, Brazil}

\AtBeginDocument{
    \providecommand\BibTeX{{Bib\TeX}}
}

%% Only for papers written in PORTUGUESE:
%% comment out the following three lines
\renewcommand\abstractname{Resumo}
\renewcommand\keywordsname{Palavras-chave}
\renewcommand\refname{Referências}
%% Rótulos de figuras e tabelas em português.
\AtBeginDocument{\renewcommand\figurename{Figura}\renewcommand\tablename{Tabela}}

%% Hifenização em português.
\usepackage[brazilian]{babel}
%% Gráfico de dispersão (Figura 1) e colunas com quebra de linha nas tabelas.
\usepackage{pgfplots}
\pgfplotsset{compat=1.18}
\usetikzlibrary{babel}
\usepackage{array}
\raggedbottom

\begin{document}

\title{Como a Comunidade Testa Sistemas RAG e Agentes de LLM? Evidências da Literatura Cinza e de Repositórios de Código}

\author{Icaro Santos de Oliveira}
\affiliation{
  \institution{Universidade Estadual do Ceará (UECE)}
  \city{Fortaleza}
  \state{CE}
  \country{Brazil}
}
\email{icaro.santos@aluno.uece.br}

\author{João Matheus de Lima Alves}
\affiliation{
  \institution{Universidade Estadual do Ceará (UECE)}
  \city{Fortaleza}
  \state{CE}
  \country{Brazil}
}
\email{joaomatheus.lima@aluno.uece.br}

\author{Ismayle de Sousa Santos}
\affiliation{
  \institution{Universidade Estadual do Ceará (UECE)}
  \city{Fortaleza}
  \state{CE}
  \country{Brazil}
}
\email{ismayle.santos@uece.br}

\renewcommand{\shortauthors}{Oliveira et al.}
\renewcommand{\shorttitle}{Como a Comunidade Testa Sistemas RAG e Agentes de LLM}

\begin{abstract}
Sistemas baseados em grandes modelos de linguagem (\textit{Large Language Models} -- LLM) com geração aumentada por recuperação (\textit{Retrieval-Augmented Generation} -- RAG) e agentes autônomos chegaram à produção sem práticas consolidadas de verificação. Este artigo investiga quais ferramentas, métodos e oráculos praticantes adotam e se o engajamento comunitário reflete a adoção técnica. Conduziu-se um estudo de métodos mistos sobre 16.836 itens únicos de cinco fontes via interfaces oficiais, 6.737 repositórios enriquecidos e 13.923 detecções em 10.123 unidades, classificando ferramentas em menção (N1), uso declarado (N2) e integração contínua (N3) sem \textit{frameworks} do sistema sob teste e repositórios mantenedores. Discurso e prática externa exibem convergência moderada ($\tau_b = 0{,}54$ em N1 $\times$ N3; $0{,}56$ em N2 $\times$ N3), sem sustentar P1: Ragas lidera menções (145) e uso (330), enquanto \textit{promptfoo} (47), LangSmith (45), SWE-bench (41) e Langfuse (40) lideram em N3. LLM como juiz lidera os métodos (671); oráculos de referência (56,5\%) superam pseudoautomáticos (52,0\%), e derivados seguem raros (4,3\%). Sobre 7.339 itens elegíveis, adoção e engajamento não se associam ($\rho = -0{,}21$, sustentando P2), e agentes superam RAG no modelo ajustado ($\beta = 0{,}023$; $p = 0{,}044$).
\end{abstract}

\keywords{Teste de Sistemas Inteligentes, Geração Aumentada por Recuperação, Agentes de LLM, Literatura Cinza, Revisão Multivocal, Mineração de Repositórios de Software}

\frenchspacing
\maketitle

% ============================================================
\section{Introdução}
\label{sec:intro}
% ============================================================

Sistemas baseados em grandes modelos de linguagem (\textit{Large Language Models} -- LLM) passaram a combinar geração aumentada por recuperação (\textit{Retrieval-Augmented Generation} -- RAG) com agentes que planejam e acionam ferramentas externas~\cite{lewis2020rag,gao2023ragsurvey,yehudai2025agents}. Essa arquitetura traz falhas não determinísticas, deriva de contexto e erros em múltiplas etapas, exigindo verificação contínua além de testes unitários tradicionais~\cite{zhang2022mltesting,barnett2024seven,cemri2025multiagent}.

Levantamentos acadêmicos recentes sistematizam métricas, conjuntos de avaliação e taxonomias de defeitos para LLMs, RAG e agentes~\cite{chang2024survey,yu2024ragsurvey,yehudai2025agents}. Esses levantamentos, contudo, mapeiam a fronteira da pesquisa publicada e não distinguem ferramentas populares em tutoriais daquelas sustentadas em fluxos de integração contínua (\textit{Continuous Integration} -- CI) em produção~\cite{borges2018star,raulamo2017choosing}.

Para preencher essa lacuna, este artigo investiga como a comunidade de engenharia de software testa sistemas RAG e agentes na prática. Adota-se uma revisão multivocal que combina literatura cinza e mineração de repositórios~\cite{garousi2019guidelines,kalliamvakou2014promises}, estruturada em cinco questões de pesquisa (QPs) sobre ferramentas, métodos, oráculos, engajamento, falhas e transferência entre paradigmas.

Em relação à coleta anterior, o \textit{corpus} foi ampliado para 16.836 itens únicos de cinco fontes coletadas via interfaces de programação de aplicações (\textit{Application Programming Interfaces} -- APIs) oficiais, com 6.737 repositórios enriquecidos e 13.923 detecções. A divisão em três níveis de evidência preserva o contraste entre discurso e CI, mantendo ferramentas de teste separadas do sistema sob teste (\textit{System Under Test} -- SUT).

As contribuições incluem um \textit{corpus} rastreável de 16.836 itens em 10.123 unidades, uma taxonomia empírica de 33 ferramentas e 10 métodos alinhados à escada de oráculos de Barr et al.~\cite{barr2015oracle}, o cálculo completo de engajamento sobre 7.339 itens elegíveis com separação dos repositórios mantenedores, e a revisão de teses iniciais sobre rastreamento e oráculos. Apenas a codificação qualitativa manual aguarda conclusão.

% ============================================================
\section{Fundamentação}
\label{sec:fundamentacao}
% ============================================================

Esta seção apresenta os conceitos que fundamentam a investigação empírica do estudo. A Seção~\ref{sec:fund-rag-agentes} caracteriza as arquiteturas de sistemas RAG e de agentes baseados em LLMs quanto às suas superfícies de falha. Em seguida, a Seção~\ref{sec:fund-oraculo} discute o problema do oráculo em saídas não determinísticas, e a Seção~\ref{sec:fund-cinza} distingue popularidade discursiva de adoção técnica em repositórios.

\subsection{Sistemas RAG e Agentes de LLM}
\label{sec:fund-rag-agentes}

Um sistema RAG acopla um módulo de recuperação de informação (\textit{Information Retrieval} -- IR) a um modelo gerador para ancorar respostas em documentos externos~\cite{lewis2020rag,gao2023ragsurvey}. As falhas ocorrem na segmentação, na busca semântica, na ordenação de passagens ou na síntese infiel ao contexto recuperado~\cite{barnett2024seven,yu2024ragsurvey}.

Agentes baseados em LLMs estendem esse ciclo ao manter estado, decompor objetivos, invocar APIs e avaliar resultados intermediários em laços iterativos~\cite{yehudai2025agents,liu2024agentbench}. A execução em múltiplas voltas introduz falhas de planejamento, chamadas inválidas de ferramentas e perda de alinhamento entre agentes~\cite{cemri2025multiagent,yao2024taubench}.

A composição híbrida, na qual agentes consultam bases vetoriais como ferramentas de busca, torna a fronteira de teste permeável entre os dois paradigmas~\cite{gao2023ragsurvey,yehudai2025agents}. Essa convergência arquitetural motiva avaliar em conjunto ferramentas de avaliação de recuperação, qualidade textual e trajetória de decisão~\cite{es2024ragas,jimenez2024swebench}.

\subsection{O Problema do Oráculo em Sistemas Baseados em LLM}
\label{sec:fund-oraculo}

Na definição clássica de Barr et al.~\cite{barr2015oracle}, um oráculo de teste distingue comportamentos aceitáveis de falhos para cada entrada executada. Em sistemas tradicionais, asserções determinísticas verificam contratos exatos de saída, mas respostas em linguagem natural admitem múltiplas formulações corretas~\cite{zhang2022mltesting,chang2024survey}.

Para contornar saídas abertas, métodos recentes empregam modelos avaliadores como juízes automáticos de coerência, fidelidade e utilidade~\cite{zheng2023judging,es2024ragas}. Essa alternativa sofre de vieses de posição, verbosidade e autoconsistência, exigindo calibração contra julgamento humano ou referências curadas~\cite{zheng2023judging,chang2024survey}.

Métodos derivados, como testes metamórficos e baseados em propriedades, dispensam respostas exatas ao verificar invariantes entre perturbações controladas da entrada~\cite{barr2015oracle,zhang2022mltesting}. Neste estudo, cada método extraído é mapeado em cinco degraus de automação do oráculo: especificado, referência, pseudoautomático, derivado e humano~\cite{barr2015oracle,zheng2023judging}.

\subsection{Literatura Cinza e Popularidade versus Adoção}
\label{sec:fund-cinza}

Revisões multivocais combinam artigos revisados por pares e literatura cinza produzida por praticantes em repositórios, fóruns e registros técnicos~\cite{garousi2019guidelines,cruzes2011thematic}. Essa triangulação captura ferramentas emergentes que ainda não aparecem em conferências acadêmicas~\cite{garousi2019guidelines,raulamo2017choosing}.

Contudo, a visibilidade de uma ferramenta em discussões ou estrelas de repositório não garante seu uso efetivo em projetos reais~\cite{borges2018star,kalliamvakou2014promises}. Repositórios de exemplo, \textit{forks} e guias introdutórios inflam contagens brutas caso não haja separação entre menção superficial e configuração executável~\cite{kalliamvakou2014promises,ralph2020standards}.

Por essa razão, a análise adota o princípio metodológico aplicado por Raulamo-Jurvanen et al.~\cite{raulamo2017choosing} e distingue três níveis graduais de evidência. O contraste entre debate discursivo e configuração verificável em CI testa empiricamente até que ponto popularidade reflete adoção sustentada~\cite{borges2018star,kalliamvakou2014promises}.

% ============================================================
\section{Método}
\label{sec:metodo}
% ============================================================

O desenho empírico segue um protocolo pré-registrado de métodos mistos que integra literatura cinza e mineração de repositórios~\cite{msrkit2027}. A Seção~\ref{sec:met-qps} define as questões de pesquisa e proposições, enquanto a Seção~\ref{sec:met-coleta} descreve a coleta automatizada estruturada pelo fluxo PRISMA 2020~\cite{page2021prisma}. Na sequência, as Seções~\ref{sec:met-niveis},~\ref{sec:met-descoberta} e~\ref{sec:met-estatistica} detalham os três níveis de evidência, a descoberta aberta de termos e o plano de análise estatística.

\subsection{Questões de Pesquisa e Proposições}
\label{sec:met-qps}

O estudo adota desenho de métodos mistos para responder a cinco questões de pesquisa definidas no protocolo pré-registrado no repositório do instrumento de coleta~\cite{msrkit2027,ralph2020standards}. Cada questão articula evidências quantitativas de mineração com categorias temáticas para triangulação~\cite{cruzes2011thematic,garousi2019guidelines}.

A Tabela~\ref{tab:qps} resume as cinco questões de pesquisa, os construtos operacionais e as fontes de evidência associadas. Todas as análises descritivas e testes não paramétricos seguem o protocolo publicado junto ao conjunto de dados~\cite{msrkit2027,page2021prisma}.

\begin{table}[htbp]
\centering
\caption{Questões de pesquisa, construtos e fontes.}
\label{tab:qps}
\footnotesize
\setlength{\tabcolsep}{4pt}
\begin{tabular}{@{}p{0.8cm}p{3.6cm}p{3.4cm}@{}}
\toprule
\textbf{QP} & \textbf{Pergunta / Construto} & \textbf{Dados e Análise} \\
\midrule
QP1 & Ferramentas usadas e diferença entre menção (N1), uso (N2) e CI (N3) & Matriz ferramenta $\times$ nível; correlação $\tau_b$ de Kendall \\
QP2 & Métodos de teste e degraus de oráculo empregados na prática & Contagens de 10 métodos e 5 degraus; coocorrência \textit{Lift}/Jaccard \\
QP3 & Relação entre práticas de teste e métricas de engajamento & Kruskal--Wallis, Dunn--Holm, Spearman, Cliff e OLS (7.339 itens) \\
QP4 & Falhas relatadas pelos praticantes na depuração e avaliação & Codificação temática na amostra triada (pendente) \\
QP5 & Sobreposição de ferramentas e métodos entre RAG e agentes & Índice de Jaccard e razão de transferência \\
\bottomrule
\end{tabular}
\end{table}

As questões quantitativas avaliam proposições orientadoras do protocolo. A proposição P1 postula divergência entre menções e adoção externa em CI ($\tau_b < 0{,}5$)~\cite{kendall1938,borges2018star}. Já P2 prevê associação fraca entre adoção e engajamento ($\rho < 0{,}4$)~\cite{kalliamvakou2014promises}, P3 contrasta agentes e RAG em 2025--2026, e a hipótese H2 avalia atomicidade de oráculos na QP2 após codificação manual.

\subsection{Coleta de Dados e Protocolo PRISMA}
\label{sec:met-coleta}

A coleta seguiu as diretrizes de fluxo do protocolo PRISMA (\textit{Preferred Reporting Items for Systematic Reviews and Meta-Analyses})~\cite{page2021prisma}. Todas as requisições foram executadas em 06/11/2026 (01h19--03h07 em Tempo Universal Coordenado -- UTC) na execução identificada por \texttt{2026-11-06T01-19-19Z-015e} pelo minerador \textit{MSR-Kit}~\cite{msrkit2027}.

O acesso respeitou exclusivamente APIs oficiais e limites explícitos de taxa, abrangendo a janela de janeiro de 2023 a agosto de 2026~\cite{msrkit2027,ralph2020standards}. Cinco fontes integraram a execução: GitHub\footnote{\url{https://github.com}}, Hacker News\footnote{\url{https://news.ycombinator.com}}, Stack Exchange\footnote{\url{https://stackexchange.com}}, Hugging Face\footnote{\url{https://huggingface.co}} e dev.to\footnote{\url{https://dev.to}}.

Fontes adicionais previstas em versões anteriores, como Reddit\footnote{\url{https://www.reddit.com}} e \textit{feeds} públicos, não compuseram esta captura, enquanto Bluesky\footnote{\url{https://bsky.app}}, Discord\footnote{\url{https://discord.com}}, LinkedIn\footnote{\url{https://www.linkedin.com}} e X\footnote{\url{https://x.com}} permaneceram indisponíveis ou restritos por política de acesso~\cite{msrkit2027}. Foram executadas 565 consultas de 20 termos em inglês e 17 grupos de consultas literais no GitHub.

Com exceção do dev.to, varrido em censo completo, aplicou-se estratificação anual com limite de 100 resultados por consulta~\cite{msrkit2027}. Ao todo, 2.771 requisições \textit{Hypertext Transfer Protocol} (HTTP) identificaram 78.308 registros brutos, dos quais 58.716 foram descartados no filtro lexical e 2.756 na remoção de duplicatas de \textit{Uniform Resource Locator} (URL), conteúdo e similaridade.

A Tabela~\ref{tab:coleta} detalha o volume por plataforma nos 16.836 itens únicos preservados. O repositório de código responde por 79,8\% do \textit{corpus}, embora 122 consultas tenham atingido o teto de 100 itens por estrato, concentrando-se em trechos de código (4.437), \textit{issues} (7.339) e metadados de repositórios (1.666).

\begin{table*}[tb]
\centering
\caption{Fontes coletadas, volume bruto, consultas e itens únicos na execução 015e.}
\label{tab:coleta}
\small
\begin{tabular}{@{}lrrrrl@{}}
\toprule
\textbf{Fonte} & \textbf{Consultas} & \textbf{Requisições} & \textbf{Bruto} & \textbf{Únicos} & \textbf{Observações da captura} \\
\midrule
GitHub & 384 & 1.017 & 15.053 & 13.442 & 7.339 \textit{issues}, 4.437 código, 1.666 repositórios; 104 truncadas \\
Hacker News & 80 & 160 & 2.455 & 1.690 & Discussões técnicas e lançamentos; 13 truncadas \\
Stack Exchange & 80 & 885 & 996 & 813 & Perguntas e respostas técnicas; 3 truncadas \\
Hugging Face & 20 & 59 & 580 & 568 & Fóruns e modelos; 2 truncadas; 38 fora da janela \\
dev.to & 1 & 650 & 508 & 323 & Censo único com 58.716 itens sem correspondência \\
\midrule
\textbf{Total} & \textbf{565} & \textbf{2.771} & \textbf{19.592*} & \textbf{16.836} & \textbf{122 consultas truncadas; 2.756 duplicatas removidas} \\
\bottomrule
\end{tabular}
\par\vspace{2pt}{\footnotesize *Após descarte de 58.716 registros do dev.to no filtro de termos sobre os 78.308 itens identificados inicialmente.\par}
\end{table*}

\subsection{Enriquecimento e Níveis de Evidência}
\label{sec:met-niveis}

Os itens vinculados ao GitHub acionaram o enriquecimento automático de repositórios para coletar sinais estruturais de verificação~\cite{kalliamvakou2014promises,msrkit2027}. Ao todo, 6.737 repositórios distintos foram enriquecidos, revelando 4.366 projetos com diretórios de CI e 5.350 com sinais ou suítes explícitas de testes.

Um catálogo estruturado definiu 51 entradas iniciais: 33 ferramentas de teste, 8 \textit{frameworks} de orquestração classificados como SUT e 10 métodos de teste~\cite{msrkit2027}. O alinhamento lexical consolidou 13.923 detecções brutas (10.748 de ferramentas e 3.175 de métodos) por 8.265 itens em 10.123 unidades, das quais 6.729 vinculam-se ao GitHub.

Para evitar autoenviesamento na contagem, o roteiro analítico remove ocorrências nos repositórios mantenedores das próprias ferramentas, retendo apenas adoção externa (\textit{N3 externo})~\cite{kalliamvakou2014promises,msrkit2027}. Cada ocorrência restante recebe classificação em três níveis graduais de evidência~\cite{raulamo2017choosing,borges2018star}.

O nível N1 (\textbf{menção}) registra citação discursiva sem vínculo a código executável, e N2 (\textbf{uso declarado}) captura dependências ou importações. Já N3 (\textbf{adoção em CI}) exige presença do identificador em fluxos de automação contínua externos ao projeto mantenedor.

\subsection{Descoberta Aberta e Saturação}
\label{sec:met-descoberta}

Para mitigar o viés de catálogo fechado, o módulo de descoberta aplicou representação por frequência do termo pelo inverso da frequência nos documentos (\textit{Term Frequency--Inverse Document Frequency} -- TF-IDF) e Alocação Latente de Dirichlet (\textit{Latent Dirichlet Allocation} -- LDA) sobre o \textit{corpus}~\cite{blei2003lda,cruzes2011thematic}. Documentos muito curtos foram excluídos antes da vetorização, restando 10.438 textos e 188.212 termos.

A triagem gerou 114 termos candidatos para inspeção manual, mantendo zero promoções automáticas para preservar a validade do catálogo~\cite{msrkit2027}. Paralelamente, ordenou-se o \textit{corpus} em 674 lotes cronológicos para avaliar o ponto de saturação da taxonomia de códigos~\cite{cruzes2011thematic,ralph2020standards}.

\subsection{Plano Estatístico e Elegibilidade de Engajamento}
\label{sec:met-estatistica}

A concordância entre rankings nos níveis N1, N2 e N3 externo usa o coeficiente $\tau_b$ de Kendall com intervalos de confiança de 95\% via 10.000 reamostragens \textit{bootstrap}~\cite{kendall1938,efron1993bootstrap}. Associações entre pares empregam o índice de Jaccard e a razão de elevação (\textit{Lift}).

Na QP3, aplica-se um filtro de elegibilidade sobre os 16.836 itens para isolar discussões maduras em plataformas com suporte a reações. Excluem-se 6.103 registros sem métrica de reação (código e repositórios), 2.505 \textit{issues} com padrão de dor ou falha, 839 itens recentes (menores de três meses) e 50 postagens automatizadas, restando 7.339 itens válidos~\cite{msrkit2027,kalliamvakou2014promises}.

Sobre os 7.339 itens válidos, calcula-se o engajamento $E \in [0,1]$ por posto percentual intraestrato. Ferramentas e métodos com $n \ge 10$ são testados via Kruskal--Wallis, Dunn--Holm, Spearman com adoção externa, binomial negativa, delta de Cliff e mínimos quadrados ordinários (\textit{Ordinary Least Squares} -- OLS) com erros robustos~\cite{kruskal1952,dunn1964,holm1979,cliff1993dominance}.

Apenas a triagem manual em duas trilhas e a dupla codificação dos 16.836 itens permanecem pendentes. Essa etapa futura responderá à QP4, verificará a hipótese H2 e resolverá a QP5 com concordância kappa de Cohen~\cite{cohen1960kappa,landis1977measurement}.

% ============================================================
\section{Resultados}
\label{sec:resultados}
% ============================================================

Esta seção apresenta os resultados empíricos sobre os 16.836 itens e as 10.123 unidades analíticas da coleta oficial. A Seção~\ref{sec:res-qp1} examina a adoção externa das 33 ferramentas de teste nos três níveis de evidência (QP1), a Seção~\ref{sec:res-qp2} analisa os dez métodos e a escada de oráculos (QP2), e a Seção~\ref{sec:res-descoberta} reporta a descoberta aberta e a saturação do catálogo. Em seguida, a Seção~\ref{sec:res-qp3} avalia a relação entre engajamento e adoção (QP3), e a Seção~\ref{sec:res-qp4-qp5} sintetiza o estado da triagem qualitativa e a sobreposição entre paradigmas (QP4 e QP5).

\subsection{QP1: Ferramentas de Teste e Níveis de Evidência}
\label{sec:res-qp1}

Após remover 27 ocorrências em repositórios mantenedores (\textit{self}) de 27 ferramentas, as 33 ferramentas de teste somam 3.538 ocorrências externas: 855 em N1, 2.018 em N2 e 665 em N3 externo distribuídas por 32 ferramentas. Os 8 \textit{frameworks} de SUT ocorrem em 29\% das unidades, liderados por LangChain (1.081) e LangGraph (635), fora dos rankings principais.

A Tabela~\ref{tab:qp1} lista as 24 ferramentas líderes em adoção externa contínua (N3), a síntese das nove restantes e os totais por nível. A liderança em N3 pertence ao \textit{promptfoo} (47)~\cite{promptfoo}, seguido por LangSmith (45), SWE-bench (41)~\cite{jimenez2024swebench}, Langfuse (40) e lm-evaluation-harness (39).

\begin{table*}[tb]
\centering
\caption{Distribuição externa das 33 ferramentas de teste nos três níveis de evidência (excluídos \textit{self} e 8 \textit{frameworks} de SUT).}
\label{tab:qp1}
\footnotesize
\begin{tabular}{@{}rllrrrrrr@{}}
\toprule
\textbf{\#} & \textbf{Ferramenta} & \textbf{Família} & \textbf{N3 (CI)} & \textbf{N2 (Uso)} & \textbf{N1 (Menção)} & \textbf{Total} & \textbf{RAG} & \textbf{Agente} \\
\midrule
1 & promptfoo~\cite{promptfoo} & avaliação LLM/RAG & 47 & 125 & 53 & 225 & 26 & 48 \\
2 & LangSmith & rastreamento & 45 & 266 & 109 & 420 & 130 & 156 \\
3 & SWE-bench~\cite{jimenez2024swebench} & agentes & 41 & 38 & 69 & 148 & 35 & 83 \\
4 & Langfuse & rastreamento & 40 & 213 & 61 & 314 & 65 & 113 \\
5 & lm-evaluation-harness & avaliação LLM/RAG & 39 & 101 & 76 & 216 & 20 & 29 \\
6 & Ragas~\cite{es2024ragas} & avaliação LLM/RAG & 36 & 330 & 145 & 511 & 210 & 106 \\
7 & DeepEval~\cite{deepeval} & avaliação LLM/RAG & 34 & 208 & 83 & 325 & 88 & 93 \\
8 & WebArena~\cite{zhou2024webarena} & agentes & 33 & 13 & 19 & 65 & 9 & 20 \\
9 & BFCL & agentes & 32 & 41 & 22 & 95 & 8 & 18 \\
10 & BEIR & recuperação & 30 & 61 & 8 & 99 & 18 & 7 \\
11 & garak~\cite{derczynski2024garak} & guardrails/segurança & 30 & 45 & 20 & 95 & 12 & 18 \\
12 & Guardrails AI & guardrails/segurança & 29 & 17 & 6 & 52 & 5 & 6 \\
13 & Inspect AI & avaliação LLM/RAG & 28 & 71 & 12 & 111 & 5 & 24 \\
14 & MTEB & recuperação & 26 & 35 & 7 & 68 & 9 & 5 \\
15 & $\tau$-bench~\cite{yao2024taubench} & agentes & 22 & 6 & 11 & 39 & 6 & 15 \\
16 & Giskard & avaliação LLM/RAG & 20 & 15 & 13 & 48 & 6 & 8 \\
17 & PyRIT & guardrails/segurança & 19 & 53 & 11 & 83 & 8 & 11 \\
18 & Opik & rastreamento & 19 & 35 & 8 & 62 & 11 & 16 \\
19 & ranx & recuperação & 17 & 51 & 0 & 68 & 2 & 0 \\
20 & AgentDojo~\cite{debenedetti2024agentdojo} & agentes & 15 & 27 & 5 & 47 & 4 & 14 \\
21 & AgentBench~\cite{liu2024agentbench} & agentes & 12 & 41 & 8 & 61 & 5 & 15 \\
22 & Arize Phoenix & rastreamento & 12 & 30 & 35 & 77 & 27 & 33 \\
23 & LLM Guard & guardrails/segurança & 9 & 29 & 7 & 45 & 4 & 6 \\
24 & OpenAI Evals & avaliação LLM/RAG & 7 & 76 & 15 & 98 & 17 & 36 \\
\midrule
25--33 & Demais 9 ferramentas* & 4 famílias & 23 & 91 & 52 & 166 & 56 & 49 \\
\midrule
\multicolumn{3}{@{}l}{\textbf{Total (33 ferramentas de teste)}} & \textbf{665} & \textbf{2.018} & \textbf{855} & \textbf{3.538} & \textbf{786} & \textbf{929} \\
\bottomrule
\end{tabular}
\par\vspace{2pt}{\footnotesize *NeMo Guardrails (4/28/12), autoevals (4/9/1), continuous-eval (4/6/2), agentevals (3/14/3), ir\_measures (3/5/1), TruLens (2/12/23), UpTrain (2/2/7), RAGChecker (1/3/2) e pytrec\_eval (0/12/1).\par}
\end{table*}

O Ragas~\cite{es2024ragas} mantém a liderança em volume total (511), uso declarado (330) e menções (145), seguido por LangSmith (420), DeepEval (325)~\cite{deepeval} e Langfuse (314). Por famílias, avaliação de LLM/RAG reúne 33,7\% das ocorrências em N3 externo (224), agentes somam 23,8\% (158), rastreamento 17,4\% (116), segurança 13,7\% (91) e recuperação 11,4\% (76).

Mesmo após isolar repositórios mantenedores, o rastreamento sustenta 116 ocorrências externas em CI (LangSmith 45, Langfuse 40, Opik 19 e Arize Phoenix 12). Ainda assim, P1 não se confirma: Kendall marca $\tau_b = 0{,}54$ [0,42; 0,59] entre N1 e N3 ($p < 0{,}001$), $\tau_b = 0{,}56$ [0,44; 0,60] entre N2 e N3 ($p < 0{,}001$) e $\tau_b = 0{,}52$ [0,41; 0,56] entre N1 e N2 ($p < 0{,}001$).

A Figura~\ref{fig:dispersao} posiciona as 33 ferramentas entre menções ($\text{N1}+1$) e adoção externa em CI ($\text{N3}+1$). À esquerda da diagonal predominam ferramentas com alta conversão para CI, como Guardrails AI (29 em N3 para 6 em N1), BEIR (30 para 8) e ranx (17 para 0); à direita acumulam-se soluções amplamente discutidas, como Ragas (145 para 36) e LangSmith (109 para 45).

\begin{figure}[htbp]
\centering
\begin{tikzpicture}
\begin{loglogaxis}[
    width=\columnwidth,
    height=6.8cm,
    xlabel={Menções ($\text{N1}+1$, escala logarítmica)},
    ylabel={Adoção externa em CI ($\text{N3}+1$, escala logarítmica)},
    xmin=0.8, xmax=260,
    ymin=0.8, ymax=85,
    xtick={1,2,6,11,51,151},
    xticklabels={0,1,5,10,50,150},
    ytick={1,2,6,11,51},
    yticklabels={0,1,5,10,50},
    grid=major,
    grid style={dotted,gray!45},
    tick label style={font=\scriptsize},
    label style={font=\footnotesize},
]
\addplot[dashed,gray!70,domain=0.8:70,samples=2] {x};
\addplot[only marks,mark=*,mark size=1.8pt,blue!75!black] coordinates {
    (54,48) (110,46) (70,42) (62,41) (77,40) (146,37) (84,35)
    (20,34) (23,33) (9,31) (21,31) (7,30) (13,29) (8,27)
    (12,23) (14,21) (12,20) (9,20) (1,18) (6,16) (9,13)
    (36,13) (8,10) (16,8) (13,5) (2,5) (3,5) (4,4) (2,4)
    (24,3) (8,3) (3,2) (2,1)
};
\node[font=\tiny,anchor=south east,yshift=1pt] at (axis cs:54,48) {promptfoo};
\node[font=\tiny,anchor=south west,xshift=-2pt,yshift=1pt] at (axis cs:110,46) {LangSmith};
\node[font=\tiny,anchor=south,yshift=2pt] at (axis cs:70,42) {SWE-bench};
\node[font=\tiny,anchor=west,xshift=2pt] at (axis cs:62,41) {Langfuse};
\node[font=\tiny,anchor=north,yshift=-2pt] at (axis cs:146,37) {Ragas};
\node[font=\tiny,anchor=north,yshift=-2pt] at (axis cs:84,35) {DeepEval};
\node[font=\tiny,anchor=south,yshift=2pt] at (axis cs:20,34) {WebArena};
\node[font=\tiny,anchor=east,xshift=-2pt] at (axis cs:7,30) {Guardrails AI};
\node[font=\tiny,anchor=south east,xshift=1pt,yshift=2pt] at (axis cs:9,31) {BEIR};
\node[font=\tiny,anchor=west,xshift=2pt] at (axis cs:1,18) {ranx};
\node[font=\tiny,anchor=west,xshift=2pt] at (axis cs:36,13) {Arize Phoenix};
\node[font=\tiny,anchor=west,xshift=2pt] at (axis cs:24,3) {TruLens};
\node[font=\tiny,anchor=south west,xshift=1pt] at (axis cs:2,1) {pytrec\_eval};
\end{loglogaxis}
\end{tikzpicture}
\caption{Dispersão entre menções (N1) e adoção externa em CI (N3) das 33 ferramentas; a linha tracejada indica N3 = N1.}
\label{fig:dispersao}
\Description{Gráfico de dispersão em escala logarítmica comparando menções (N1) e adoção externa em integração contínua (N3) das 33 ferramentas de teste.}
\end{figure}

O peso de fontes ancoradas em código explicita assimetrias de visibilidade: ferramentas como ranx (zero N1), AgentBench (8 em N1)~\cite{liu2024agentbench}, LLM Guard (7 em N1) e MTEB (7 em N1) concentram-se em repositórios. A Tabela~\ref{tab:pares} reúne os principais pares de coocorrência externos entre os 252 pares analisados sobre 10.123 unidades.

\begin{table}[htbp]
\centering
\caption{Pares selecionados de coocorrência externa ($N = 10.123$ unidades).}
\label{tab:pares}
\footnotesize
\setlength{\tabcolsep}{3.5pt}
\begin{tabular}{@{}llrrr@{}}
\toprule
\textbf{Grupo} & \textbf{Par} & \textbf{Unid.} & \textbf{Jaccard} & \textbf{\textit{Lift}} \\
\midrule
F $\times$ F & Guardrails AI + NeMo Guardrails & 10 & 0,12 & 44,2 \\
F $\times$ F & garak + PyRIT & 21 & 0,13 & 27,0 \\
F $\times$ F & garak + promptfoo & 34 & 0,12 & 16,1 \\
F $\times$ F & Arize Phoenix + Giskard & 5 & 0,04 & 13,7 \\
F $\times$ M & SWE-bench + teste metamórfico & 8 & 0,04 & 13,3 \\
F $\times$ M & Arize Phoenix + avaliação humana & 9 & 0,05 & 9,8 \\
F $\times$ M & garak + \textit{red teaming} & 21 & 0,07 & 9,2 \\
F $\times$ M & SWE-bench + avaliação humana & 15 & 0,06 & 8,5 \\
F $\times$ SUT & TruLens + LlamaIndex & 8 & 0,03 & 10,9 \\
F $\times$ SUT & Arize Phoenix + AutoGen & 6 & 0,04 & 9,5 \\
F $\times$ SUT & Guardrails AI + CrewAI & 5 & 0,03 & 7,6 \\
F $\times$ SUT & DeepEval + Haystack & 10 & 0,03 & 7,4 \\
\bottomrule
\end{tabular}
\par\vspace{2pt}{\footnotesize Nota: F indica ferramenta de teste; M, método; SUT, sistema sob teste.\par}
\end{table}

Os pares de maior razão de elevação na Tabela~\ref{tab:pares} revelam composições de segurança em camadas, como Guardrails AI com NeMo Guardrails (\textit{Lift} 44,2), \textit{garak} com PyRIT (\textit{Lift} 27,0) e \textit{garak} com \textit{promptfoo} (34 unidades; \textit{Lift} 16,1). Destacam-se ainda combinações específicas de orquestração e avaliação, como TruLens com LlamaIndex (\textit{Lift} 10,9) e Arize Phoenix com AutoGen (\textit{Lift} 9,5).

\subsection{QP2: Métodos e Escada de Oráculos}
\label{sec:res-qp2}

Os dez métodos do catálogo totalizam 2.269 ocorrências, das quais 1.290 unidades reúnem ao menos um sinal de oráculo dentre os sete métodos mapeados na taxonomia. Como mostra a Tabela~\ref{tab:qp2}, o método LLM como juiz lidera com folga (671 unidades)~\cite{zheng2023judging}, seguido por conjunto \textit{golden} (467), regressão de avaliação (368), \textit{red teaming} (244) e teste por asserção (203).

\begin{table*}[tb]
\centering
\caption{Métodos de teste identificados, degrau de oráculo e distribuição por paradigma.}
\label{tab:qp2}
\small
\begin{tabular}{@{}llrrr@{}}
\toprule
\textbf{Método} & \textbf{Degrau de oráculo} & \textbf{Unidades} & \textbf{RAG} & \textbf{Agente} \\
\midrule
LLM como juiz & pseudoautomático & 671 & 254 & 337 \\
conjunto \textit{golden} & referência & 467 & 243 & 252 \\
regressão de avaliação & referência & 368 & 150 & 239 \\
\textit{red teaming} & --- & 244 & 97 & 172 \\
teste por asserção & especificado & 203 & 103 & 153 \\
avaliação humana & humano & 121 & 81 & 82 \\
avaliação de trajetória & --- & 96 & 44 & 71 \\
geração sintética de testes & --- & 42 & 30 & 32 \\
teste metamórfico & derivado & 41 & 29 & 33 \\
teste baseado em propriedades & derivado & 16 & 6 & 11 \\
\bottomrule
\end{tabular}
\end{table*}

Métodos específicos de robustez formal ainda têm baixa penetração empírica: avaliação de trajetória aparece em 96 unidades, geração sintética de testes em 42, teste metamórfico em 41 e teste baseado em propriedades em apenas 16. Parcela expressiva das ocorrências envolve sistemas híbridos (como 176 registros simultâneos para LLM como juiz) ou itens sem rótulo de paradigma (256 no mesmo método).

A consolidação na escada de Barr et al.~\cite{barr2015oracle}, apresentada na Tabela~\ref{tab:oraculos}, revela uma inversão importante em relação à coleta preliminar. Os oráculos de referência passam à liderança global com 729 unidades (56,5\% das 1.290 com sinal), superando os pseudoautomáticos (671; 52,0\%), seguidos por especificados (203; 15,7\%), humanos (121; 9,4\%) e derivados (55; 4,3\%).

\begin{table}[htbp]
\centering
\caption{Escada de oráculos nas unidades com sinal ($N = 1.290$ todas; $465$ RAG; $656$ Agente).}
\label{tab:oraculos}
\footnotesize
\begin{tabular}{@{}lrrr@{}}
\toprule
\textbf{Degrau} & \textbf{Todas (\%)} & \textbf{RAG (\%)} & \textbf{Agente (\%)} \\
\midrule
Referência & 729 (56,5) & 322 (69,2) & 418 (63,7) \\
Pseudoautomático & 671 (52,0) & 254 (54,6) & 337 (51,4) \\
Especificado & 203 (15,7) & 103 (22,2) & 153 (23,3) \\
Humano & 121 (9,4) & 81 (17,4) & 82 (12,5) \\
Derivado & 55 (4,3) & 33 (7,1) & 42 (6,4) \\
\bottomrule
\end{tabular}
\par\vspace{2pt}{\footnotesize Nota: somas superam 100\% porque uma unidade pode combinar múltiplos degraus.\par}
\end{table}

Os valores das colunas excedem 100\% (somando 137,9\% no total, 170,5\% em RAG e 157,3\% em agentes) porque unidades combinam múltiplos oráculos em paralelo, sobretudo conjuntos \textit{golden} avaliados por juízes de LLM. Nota-se ainda que a avaliação humana cresce de 5,0\% na coleta piloto para 9,4\% no \textit{corpus} expandido, alcançando 17,4\% no recorte de RAG.

\subsection{Descoberta Aberta e Saturação do Catálogo}
\label{sec:res-descoberta}

O módulo de descoberta sobre 10.438 textos válidos extraiu 114 termos candidatos sem promoção automática ao catálogo semente de 51 entradas. O principal achado arquitetural é a emergência do \textit{Model Context Protocol} (MCP), concentrado em um grupo temático de 2.071 documentos ligados à padronização de ferramentas para agentes.

Por outro lado, a modelagem de tópicos evidenciou blocos recorrentes de ruído editorial e técnico no \textit{corpus} bruto. Destacam-se grupos de modelos automáticos de relato de erro (1.681 e 734 documentos), boletins automáticos de artigos recentes (286), referências sociais (131), alertas de pacotes vulneráveis (42) e rastros de documentação de ferramentas (26).

Na análise de saturação ao longo de 674 lotes sequenciais, o último código inédito (\texttt{tool:uptrain}) ingressou no lote 24, completando as 51 entradas da taxonomia. Os 650 lotes subsequentes não adicionaram códigos novos, estabilizando a cobertura em 5,38\% das 10.123 unidades, embora esse resultado reflita o esgotamento do dicionário fixo antes da revisão dos 114 candidatos.

\subsection{QP3: Maturidade de Testes e Engajamento}
\label{sec:res-qp3}

Do \textit{corpus} de 16.836 itens, 7.339 cumpriram os critérios de elegibilidade de engajamento (excluídos 6.103 sem métrica, 2.505 de dor, 839 recentes e 50 automatizados). Nesse recorte, 20 ferramentas ($n = 1.060$ ocorrências) e todos os 10 métodos atingiram $n \ge 10$ itens.

A Tabela~\ref{tab:qp3} apresenta as 15 ferramentas de maior mediana de engajamento $E$, seguidas das cinco elegíveis restantes em nota. Apenas MTEB ($0{,}73$) e Arize Phoenix ($0{,}66$) superam o piso comum de $0{,}43$, no qual empatam as demais 18 ferramentas e todos os 10 métodos (como LLM como juiz, $n = 669$, e conjunto \textit{golden}, $n = 484$).

\begin{table}[htbp]
\centering
\caption{Engajamento $E$ por ferramenta ($n \ge 10$; 15 primeiras de 20 elegíveis).}
\label{tab:qp3}
\footnotesize
\setlength{\tabcolsep}{3.5pt}
\begin{tabular}{@{}lrcr@{}}
\toprule
\textbf{Ferramenta} & \textbf{$n$} & \textbf{Mediana $E$ [IC 95\%]} & \textbf{N3 ext.} \\
\midrule
MTEB & 11 & 0,73 [0,43; 0,90] & 26 \\
Arize Phoenix & 35 & 0,66 [0,43; 0,89] & 12 \\
lm-eval-harness & 210 & 0,43 [0,43; 0,43] & 39 \\
Ragas~\cite{es2024ragas} & 159 & 0,43 [0,43; 0,43] & 36 \\
LangSmith & 126 & 0,43 [0,43; 0,43] & 45 \\
Langfuse & 99 & 0,43 [0,43; 0,43] & 40 \\
DeepEval~\cite{deepeval} & 90 & 0,43 [0,43; 0,43] & 34 \\
SWE-bench~\cite{jimenez2024swebench} & 70 & 0,43 [0,43; 0,43] & 41 \\
promptfoo~\cite{promptfoo} & 63 & 0,43 [0,43; 0,43] & 47 \\
TruLens & 28 & 0,43 [0,43; 0,43] & 2 \\
PyRIT & 27 & 0,43 [0,43; 0,43] & 19 \\
garak~\cite{derczynski2024garak} & 26 & 0,43 [0,43; 0,43] & 30 \\
Inspect AI & 18 & 0,43 [0,43; 0,43] & 28 \\
Opik & 18 & 0,43 [0,43; 0,43] & 19 \\
OpenAI Evals & 17 & 0,43 [0,43; 0,43] & 7 \\
\bottomrule
\end{tabular}
\par\vspace{2pt}{\footnotesize Nota: WebArena ($n=15$, 0,43 [0,43; 0,66]), BFCL (14), Giskard (13), NeMo Guardrails (11) e AgentBench (10) também apresentam mediana 0,43 [0,43; 0,43].\par}
\end{table}

Esse empate generalizado decorre da composição da amostra: 56,1\% dos 7.339 itens elegíveis são \textit{issues} do GitHub, das quais 86,0\% (3.544 de 4.120) possuem zero reações, fixando o posto percentual médio em $E = 0{,}4301$. O teste de Kruskal--Wallis detecta variação global discreta ($H = 34{,}86$; $p = 0{,}015$; $\varepsilon^2 = 0{,}033$), sem pares significativos após Dunn--Holm (menor $p_{\text{Holm}} = 0{,}162$).

Nos quadrantes, apenas MTEB e Arize Phoenix combinam alto engajamento e baixa adoção; dez ferramentas unem baixo engajamento com alta adoção, oito ocupam o quadrante inferior duplo e zero alcança o quadrante superior duplo. Em consonância, a correlação de Spearman marca $\rho = -0{,}21$ ($p = 0{,}374$), sustentando P2 e mantendo-se sem listas agregadoras ($\rho = -0{,}26$), sem estratos truncados ($\rho = -0{,}12$ em 2.550 itens) e na binomial negativa ($\rho = -0{,}20$; $p = 0{,}400$, sem convergência numérica).

Para P3 em 2025--2026 ($n = 2.043$: 1.608 de agentes contra 435 de RAG), o contraste não ajustado é nulo ($\delta = 0{,}04$; $p = 0{,}140$), mas o modelo OLS ajustado por idade e estrato indica leve vantagem de agentes ($\beta = 0{,}023$; IC 95\% $[0{,}001; 0{,}045]$; $p = 0{,}044$). Já as 2.505 \textit{issues} de falha concentram-se em 24 ferramentas, lideradas por Arize Phoenix (398), Langfuse (311), Ragas (252) e Guardrails AI (223).

\subsection{QP4 e QP5: Falhas e Transferência entre Paradigmas}
\label{sec:res-qp4-qp5}

A resposta qualitativa para QP4 depende da triagem em duas trilhas e da dupla codificação manual dos 16.836 itens atualmente marcados como pendentes no fluxo PRISMA~\cite{page2021prisma,cruzes2011thematic}. Quando concluída, essa codificação permitirá confrontar falhas relatadas por praticantes com taxonomias de defeitos de RAG e multiagentes~\cite{barnett2024seven,cemri2025multiagent}.

Para QP5, a prévia sobre as 33 ferramentas de teste (excluídos os SUTs) indica forte convergência: 31 ocorrem em ambos os sistemas (Jaccard de 0,94), e 100\% das ferramentas em agentes aparecem também em RAG. Apenas \texttt{ir\_measures} (3) e \texttt{ranx} (2) restringem-se a RAG, sem ferramentas exclusivas de agentes.

Todos os 10 métodos e 5 degraus de oráculo também são compartilhados entre RAG e agentes. Contudo, essa elevada interseção requer cautela interpretativa: 65,5\% das ocorrências de ferramentas (2.318 de 3.538) e 27,9\% dos métodos (632 de 2.269) ficam sem rótulo de paradigma, somados a casos híbridos.

% ============================================================
\section{Discussão}
\label{sec:discussao}
% ============================================================

A ampliação do \textit{corpus} para 16.836 itens e a remoção de ocorrências mantenedoras alteram interpretações do recorte piloto~\cite{msrkit2027}. Primeiro, ferramentas de rastreamento como LangSmith (45 em N3 externo), Langfuse (40) e Opik (19) consolidam uso contínuo externo e somam 17,4\% da adoção em CI.

Segundo, P1 não se sustenta ($\tau_b = 0{,}54$ em N1 $\times$ N3; $0{,}56$ em N2 $\times$ N3)~\cite{kendall1938,borges2018star}. Ainda assim, persiste assimetria de conversão: bibliotecas de avaliação lideram em volume, enquanto \textit{benchmarks} de agentes e guardrails convertem fração maior para CI~\cite{zhou2024webarena,derczynski2024garak}.

Terceiro, oráculos de referência (56,5\%) superam pseudoautomáticos (52,0\%), mostrando que juízes de LLM complementam conjuntos \textit{golden}~\cite{barr2015oracle,zheng2023judging}. Por fim, P2 se confirma ($\rho = -0{,}21$): engajamento discursivo e adoção em CI medem fenômenos distintos, agravados pelo piso de zero reações em 86,0\% das \textit{issues} e pela concentração de relatos de falha em observabilidade~\cite{borges2018star,yehudai2025agents}.

% ============================================================
\section{Ameaças à Validade}
\label{sec:ameacas}
% ============================================================

Na \textbf{validade interna}, embora os 27 registros em repositórios mantenedores tenham sido isolados, a triagem manual pendente mantém ruído tópico e 65,5\% das detecções externas de ferramentas sem rótulo de paradigma~\cite{page2021prisma,kalliamvakou2014promises}. Além disso, 86,0\% das \textit{issues} elegíveis no GitHub não têm reações, comprimindo medianas de engajamento em $0{,}4301$.

Na \textbf{validade de construto}, o alinhamento lexical verifica presença em CI sem medir cobertura de asserções ou bloqueios em \textit{pull requests} (PRs)~\cite{kalliamvakou2014promises,raulamo2017choosing}. O modelo binomial negativo complementar não convergiu numericamente, e os 114 termos descobertos por LDA ainda exigem revisão manual~\cite{blei2003lda,cruzes2011thematic}.

Na \textbf{validade externa}, 122 consultas atingiram o teto de 100 resultados por estrato anual e a captura restringiu-se a termos em inglês em cinco plataformas públicas~\cite{garousi2019guidelines,msrkit2027}. Os achados refletem o ecossistema aberto entre 2023 e 2026 e não se generalizam diretamente a \textit{pipelines} proprietários fechados~\cite{ralph2020standards}.

% ============================================================
\section{Trabalhos Relacionados}
\label{sec:relacionados}
% ============================================================

Levantamentos recentes sistematizam a avaliação de LLMs~\cite{chang2024survey}, sistemas RAG~\cite{yu2024ragsurvey} e agentes autônomos~\cite{yehudai2025agents} com foco em métricas acadêmicas e \textit{benchmarks} padronizados, enquanto Zhang et al.~\cite{zhang2022mltesting} organizam fundamentos do teste de sistemas de aprendizado de máquina. Este estudo complementa essas sínteses ao medir a adoção prática observável em repositórios e discussões técnicas.

Na investigação de defeitos, Barnett et al.~\cite{barnett2024seven} derivam sete pontos de falha em sistemas RAG a partir de estudos de caso, e Cemri et al.~\cite{cemri2025multiagent} propõem uma taxonomia de falhas para arquiteturas multiagentes. A etapa qualitativa prevista para QP4 confrontará diretamente essas categorias com os relatos reais coletados da comunidade~\cite{cruzes2011thematic}.

Metodologicamente, o desenho aproxima-se do trabalho de Raulamo-Jurvanen et al.~\cite{raulamo2017choosing}, que empregaram literatura cinza para mapear ferramentas preferidas na indústria, e dialoga com a análise de Borges e Valente~\cite{borges2018star} sobre limites da popularidade em repositórios. Como avanço, este artigo estratifica as ocorrências em três níveis de evidência e isola \textit{frameworks} de SUT dos instrumentos de verificação~\cite{kalliamvakou2014promises}.

% ============================================================
\section{Conclusão}
\label{sec:conclusao}
% ============================================================

Este artigo investigou como a comunidade testa sistemas RAG e agentes de LLM sobre 16.836 itens, 6.737 repositórios e 10.123 unidades. Mesmo isolando repositórios mantenedores, discurso e prática externa exibem concordância moderada ($\tau_b = 0{,}54$ em N1 $\times$ N3; $0{,}56$ em N2 $\times$ N3), não sustentando P1, com liderança externa em CI de \textit{promptfoo}, LangSmith, SWE-bench e Langfuse.

Nos métodos e oráculos, juízes de LLM lideram isoladamente, mas oráculos de referência superam os pseudoautomáticos em cobertura. Na QP3 sobre 7.339 itens elegíveis, engajamento não prediz adoção ($\rho = -0{,}21$, sustentando P2), enquanto agentes superam RAG no modelo ajustado ($p = 0{,}044$) e relatos de falha somam 2.505 registros.

Como trabalho futuro imediato, resta concluir a triagem em duas trilhas e a dupla codificação manual dos 16.836 itens para avaliar a hipótese H2 e responder qualitativamente à QP4. Essa etapa também resolverá as detecções sem rótulo de paradigma na QP5 e validará manualmente os 114 termos candidatos descobertos.

\section*{Disponibilidade de Artefatos}

O protocolo, os Registros de Decisão Arquitetural (\textit{Architectural Decision Records} -- ADRs), o catálogo de ferramentas e o roteiro de análise (\path{analysis/artigo2.py}) residem no repositório oficial\footnote{\url{https://github.com/Icar0S/Miner-Tests-Agents-Rag}} sob licença Apache-2.0~\cite{msrkit2027}. O pacote de reprodução do \textit{corpus} será depositado no Zenodo\footnote{\url{https://zenodo.org}} sob licença Creative Commons Atribuição 4.0 Internacional (CC BY 4.0) com metadados, URLs e identificadores pseudonimizados, preservando textos integrais apenas quando permitido na origem.

\section*{Agradecimentos}
% TODO: agradecimentos (bolsas, laboratório, disciplina). Verificar a política
% do evento sobre declarar o uso de IA generativa na ferramenta e na escrita.
Os agradecimentos institucionais e as bolsas de fomento serão inseridos na versão final após o processo de avaliação por pares. Na mesma oportunidade, serão incluídas as declarações exigidas pelo evento quanto ao suporte computacional.

\bibliographystyle{ACM-Reference-Format}
\bibliography{referencias}

\end{document}
